\section{Détecter les courbes aberrantes}\label{sec-detection}

\begin{objectif}
(O3) Construire un test qui signale une courbe Totem aberrante avec un taux de fausses alertes connu, nomme le
mécanisme en cause, en donne l'amplitude en unités naturelles, répare la courbe, et distingue une anomalie Totem d'une
erreur de marquage FO.
\end{objectif}

\subsection{Le principe en cinq étapes}
\begin{enumerate}
\item Calculer l'écart de portage $\delta$ et la covariance de son bruit à partir des dispersions et des nombres de
contributeurs publiés par Totem (\cref{sec-noise}).
\item Retirer un écart de repo lisse, $\beta_0+\beta_1\ln T$ : c'est un désaccord légitime entre Totem et le FO, pas une
anomalie de courbe (\cref{sec-null}).
\item Tester, une par une, les formes types de la \cref{prop-signatures} par un filtre adapté, avec un seuil commun
qui fixe le taux de fausses alertes de l'ensemble des tests (\cref{sec-test}).
\item En cas d'alerte, retenir les formes qui expliquent l'écart, donner leurs amplitudes en unités naturelles et la
courbe réparée (\cref{sec-attribution}).
\item Indépendamment, tester le niveau $\beta_0$ pour signaler une erreur de marquage FO possible
(\cref{sec-attribution}).
\end{enumerate}

\subsection{Le modèle de bruit}\label{sec-noise}
\begin{definition}[Modèle d'observation]\label{def-obs}
Le détecteur modélise l'écart de portage par
\begin{equation}\label{eq-obs}
\delta=X\beta+\sum_{g\in\mathcal S}a_g\,g+\varepsilon, \qquad\varepsilon\sim\N(0, \kappa\Gamma),
\end{equation}
où $X\in\M_{N,2}(\R)$ a pour colonnes $(T_k)_k$ et $(T_k\ln T_k)_k$, de sorte que $X\beta$ est l'écart de repo lisse
$\beta_0+\beta_1\ln T$ écrit en portage ; $\mathcal S$ est un ensemble, éventuellement vide, de formes types
$g\in\R^N$ d'amplitudes $a_g$ ; $\Gamma$ est la covariance du bruit de consensus, construite ci-dessous à partir des
statistiques Totem ; et $\kappa>0$ est une échelle unique, calibrée sur un historique sain. Sous l'hypothèse nulle,
$\mathcal S=\emptyset$.
\end{definition}

D'après le \cref{lem-logf} et la \cref{prop-gap}, le bruit de consensus en portage est $\varepsilon_k=-e^{\opt}_k-
e^{\dv}_k$, où $e^{\opt}_k$ est l'erreur du consensus $\ln F^{\tot}_k$ du panel options et
$e^{\dv}_k=\delta D^{\tot}_k/(S_0-D_k)$ celle des dividendes de consensus alloués. On prend les deux panels
indépendants ($n^{\com}=0$, le cas prudent), donc $\Gamma=\Gamma^{\opt}+\Gamma^{\dv}$. Totem publie, pour chaque
maturité, la dispersion $\sigma^F_k$ des forwards des contributeurs, d'où la dispersion relative $s_k=\sigma^F_k/F^{\tot}_k$,
et, pour chaque période, la dispersion $\sigma^\ell_p$ des niveaux de dividendes.

\begin{proposition}[Covariance du bruit]\label{prop-noise}
On décompose la dispersion relative des forwards en une part repo et une part dividendes,
$s_k^2=\sigma_q^2T_k^2+s_k'^2$, où $\sigma_q$ est la dispersion des vues de repo des contributeurs. On suppose les vues
de repo persistantes en maturité, et les vues de dividendes formées d'un facteur persistant, de poids $w\in[0,1]$, et
d'incréments indépendants. Alors
\begin{align}
\Gamma^{\opt}_{kl}&=\frac{\sigma_q^2\,T_kT_l+w\,s'_ks'_l+(1-w)\min(s'_k, s'_l)^2}{\max\bigl(n^{\opt}_k, n^{\opt}_l\bigr)},
\label{eq-gamma-opt}\\
\Gamma^{\dv}&=A_S\,\Gamma^\ell\,\tr{A_S}, \qquad(A_S)_{kp}=\frac{a_{kp}}{S_0-D_k}, \qquad
\Gamma^\ell_{pp'}=\frac{\rho_{pp'}\,\sigma^\ell_p\sigma^\ell_{p'}}{\max\bigl(n^{\dv}_p, n^{\dv}_{p'}\bigr)} ,
\label{eq-gamma-div}
\end{align}
où $\rho_{pp'}$ est la corrélation des erreurs de vue entre les périodes $p$ et $p'$.
\end{proposition}
La démonstration est en \cref{app-noise}. Le prototype prend $w=1/2$, $\rho_{pp'}=e^{-|c_p-c_{p'}|/2}$ pour des périodes
de milieux $c_p$ et $c_{p'}$ (en années), et lit $\sigma_q$ sur le court terme, où l'incertitude sur les dividendes est
négligeable : c'est la médiane de $s_k/T_k$ sur les trois premiers piliers (le minimum, biaisé vers le bas,
sous-estimerait l'incertitude sur le niveau de repo). Un terme diagonal tient compte du résidu de l'ajustement de la
parité. L'échelle $\kappa$ absorbe un biais commun de ces réglages ; c'est la \emph{forme} de $\Gamma$ qui rend le
détecteur puissant, car elle lui indique quels écarts sont bon marché (lisses, longs) et lesquels sont coûteux (courts,
localisés).

\paragraph{Les pics de dispersion sont des symptômes.} Quand une partie des contributeurs place un dividende de l'autre
côté d'une échéance (\cref{prop-signatures}, (iii)), la dispersion publiée à cette échéance augmente. La garder telle
quelle dans $\Gamma$ masquerait l'anomalie même qu'elle signale.

\begin{proposition}[Masquage]\label{prop-masking}
Dans la situation (iii) de la \cref{prop-signatures}, soit $s_k$ la dispersion hors désaccord et
$\Delta=\tilde d_i/(S_0-D_k)$. La statistique du déplacement en $T_k$, calculée avec la dispersion publiée, vaut
\[
z=\frac{f\,\Delta}{\sqrt{s_k^2+f(1-f)\Delta^2}\big/\sqrt{n^{\opt}_k}}\;\le\;\sqrt{\frac{n^{\opt}_k\,f}{1-f}}
\]
quelle que soit la taille du dividende, alors qu'avec la dispersion hors désaccord, $z$ croît linéairement avec
$\Delta$.
\end{proposition}
\begin{proof}
Le déplacement moyen et la variance supplémentaire sont donnés par la \cref{prop-signatures}, (iii) ; quand
$\Delta\to+\infty$, le quotient croît vers $\sqrt{n^{\opt}_k}\,f/\sqrt{f(1-f)}$.
\end{proof}
Avec $f=\num{0.375}$ et $n^{\opt}_k=16$, la borne vaut $\num{3.1}$ : à peine le seuil du test, si bien que le désaccord
passerait souvent inaperçu. Le prototype plafonne donc un pic isolé de dispersion à $\num{1.5}$ fois la plus grande des
deux dispersions voisines avant de construire $\Gamma$, et publie le rapport à la voisine, qui est lui-même un
indicateur de désaccord. Dans l'exemple fil rouge, ce rapport vaut \vRunTimingRatio{} à six mois dans le scénario de
désaccord de date, et reste inférieur à $2$ partout ailleurs.

\subsection{Le modèle nul : ce qui n'est pas une anomalie}\label{sec-null}
Le terme $X\beta$ représente un écart de repo lisse entre Totem et le FO : $\beta_0$ est l'écart de repo à un an (car
$\ln1=0$) et $\beta_1$ sa pente par unité de $\ln T$. Un tel écart peut venir d'une vraie différence de vue sur le repo,
d'une erreur de marquage du FO ou d'un biais lisse du consensus ; dans tous les cas, ce n'est pas une courbe
\emph{aberrante}. Le détecteur le retire de chaque test et le publie à part : $\hat\beta_0$ alimente le test d'erreur de
marquage FO (\cref{sec-attribution}). Si deux paramètres se révèlent trop rigides sur les données réelles, une spline
pénalisée peut remplacer $X\beta$~\cite{eilers1996} ; rien d'autre ne change.

\subsection{Formes types et test}\label{sec-test}
Chaque mécanisme de la \cref{prop-signatures} définit une forme type $g\in\R^N$, testée à chaque position possible
(\cref{tab-templates}). Avec onze piliers, on obtient \vNTemplates{} formes types.

\begin{table}[ht]
\centering
\caption{Formes types du détecteur. $\one$ est le vecteur dont toutes les composantes valent $1$, $e_k$ le $k$-ième
vecteur de la base canonique et $K_h$ un noyau gaussien de largeur $h=20$ jours. Les unités naturelles découlent de la
\cref{prop-signatures}. Un pic au dernier pilier est le même vecteur qu'une marche à ce pilier.}
\label{tab-templates}
\small
\begin{tabular}{@{}>{\raggedright\arraybackslash}p{1.9cm}>{\raggedright\arraybackslash}p{3.6cm}
>{\raggedright\arraybackslash}p{4.2cm}>{\raggedright\arraybackslash}p{4.0cm}@{}}
\toprule
Forme type & Vecteur $g$ & Mécanismes & Amplitude $a$ en unités naturelles\\
\midrule
niveau & $\one$ & décalage du spot ; dividende avant $T_1$ & spot : $\varepsilon\approx-a$ ;
dividende : $\delta d\approx-a\,(S_0-D_1)$\\
marche en $T_k$ & $\bigl(\one_{\{l\ge k\}}\bigr)_{l\in\ient1N}$ & dividendes détachés dans $\intoc{T_{k-1}}{T_k}$ ;
sortie de contributeurs en $T_k$ & $\delta\ell\approx-a\,(S_0-D_k)$\\
pic en $T_k$ & $e_k$ & forward aberrant ; désaccord de date autour de $T_k$ & $\mu\approx-a$ ; part du dividende
$f\approx a\,(S_0-D_k)/\tilde d_i$\\
saison & $\Bigl(\sum_i\dfrac{\tilde d_i\,K_h(T_k-t_i)}{S_0-D_k}\Bigr)_k$ & profil de détachement décalé & décalage
$\Delta\approx a$ années\\
\bottomrule
\end{tabular}
\end{table}

\paragraph{Blanchiment et projection.} Soit $W\in\M_N(\R)$ une matrice de blanchiment, telle que
$W(\kappa\Gamma)\tr{W}=I_N$ (par exemple $W=L^{-1}$ avec la factorisation de Cholesky $\kappa\Gamma=L\,\tr{L}$), et
$\Pi$ la projection orthogonale sur l'orthogonal des colonnes de $WX$. On pose
\[
\eh=\Pi W\delta\quad\text{(résidu blanchi)}, \qquad\gt=\Pi Wg\quad\text{(forme type blanchie)}.
\]
Après blanchiment, le bruit est un vecteur gaussien de covariance identité, et la projection retire l'écart de repo
lisse.

\begin{theorem}[Test du rapport de vraisemblance pour une forme type]\label{thm-glrt}
Dans le modèle $\delta=X\beta+a\,g+\varepsilon$, $\varepsilon\sim\N(0, \kappa\Gamma)$, avec $\beta$ inconnu, la
statistique du rapport de vraisemblance généralisé de $H_0:a=0$ contre $H_1:a\ne0$ est $z_g^2$, avec
\begin{equation}\label{eq-z}
z_g=\frac{\sprod{\gt}{\eh}}{\norm{\gt}}\sim\N\bigl(a\norm{\gt}, 1\bigr).
\end{equation}
L'estimateur du maximum de vraisemblance de l'amplitude est $\hat a_g=\sprod{\gt}{\eh}/\norm{\gt}^2$, d'écart-type
$1/\norm{\gt}$. Parmi les tests invariants par le sous-espace du modèle nul et par le signe de $a$, celui qui rejette
pour les grandes valeurs de $|z_g|$ est uniformément le plus puissant.
\end{theorem}
\begin{proof}
On a $W\delta=WX\beta+a\,Wg+W\varepsilon$ avec $W\varepsilon\sim\N(0, I_N)$ ; la log-vraisemblance maximisée vaut, à une
constante près, $-\tfrac12$ fois la somme des carrés des résidus. Sous $H_0$, cette somme vaut $\norm{\eh}^2$ ; sous
$H_1$, d'après le théorème de Frisch, Waugh et Lovell~\cite{lovell1963}, elle vaut
$\norm{\eh}^2-\sprod{\gt}{\eh}^2/\norm{\gt}^2$. Donc $2\ln\Lambda=z_g^2$, et $\hat a_g$ est le coefficient de la
régression partielle. Comme $\Pi WX=0$ et que $\Pi$ est symétrique et idempotente,
$\eh=a\,\gt+\Pi W\varepsilon$ et $\sprod{\gt}{\eh}=a\norm{\gt}^2+\sprod{\gt}{W\varepsilon}$, où
$\sprod{\gt}{W\varepsilon}\sim\N(0, \norm{\gt}^2)$. L'optimalité est le résultat des détecteurs de sous-espace adapté de
Scharf et Friedlander~\cite{scharf1994}.
\end{proof}

\begin{corollary}[Puissance]\label{cor-power}
Au seuil $c$, une anomalie de forme $g$ est détectée avec probabilité $\Phi(\lambda-c)+\Phi(-\lambda-c)$, où $\Phi$ est
la fonction de répartition de la loi normale centrée réduite et $\lambda=|a|\,\norm{\gt}$ le \emph{rapport signal sur
bruit} de l'anomalie. Au seuil typique du banc d'essai, $c\approx\vZcrit$, une anomalie de rapport $\lambda=4$ est
détectée avec probabilité \vPowThFour{} et une anomalie de rapport $\lambda=6$ avec probabilité \vPowThSix.
\end{corollary}

\begin{proposition}[Test global et seuil commun]\label{prop-family}
Sous $H_0$ :
\begin{enumerate}[label=(\roman*)]
\item $\norm{\eh}^2$ suit une loi du $\chi^2$ à $N-2$ degrés de liberté ; ce test global détecte une rugosité
qu'aucune forme type n'explique ;
\item le vecteur $(z_g)_{g}$ est gaussien centré, de matrice de corrélation
$\Omega_{gh}=\sprod{\gt}{\tilde h}/(\norm{\gt}\norm{\tilde h})$, qui ne dépend ni de $\beta$ ni de $\kappa$ ; le seuil
$c_\alpha$ tel que $\mathbb P(\max_g|z_g|>c_\alpha)=\alpha/2$ se calcule exactement en simulant $\N(0, \Omega)$, et il
est inférieur au seuil de Šidák~\cite{sidak1967} ;
\item la règle \og alerte si $\max_g|z_g|>c_\alpha$ ou si $\norm{\eh}^2$ dépasse son quantile d'ordre $1-\alpha/2$
\fg{} a un taux de fausses alertes au plus égal à $\alpha$.
\end{enumerate}
\end{proposition}
\begin{proof}
(i) $\eh=\Pi W\varepsilon$ et $\Pi$ est une projection orthogonale de rang $N-2$. (ii) $z_g$ est une forme linéaire
de $W\varepsilon$ ; l'inégalité de Šidák donne $\mathbb P(\max_g|z_g|\le c)\ge\prod_g\mathbb P(|z_g|\le c)$ pour toute
corrélation. (iii) Inégalité de la réunion.
\end{proof}

\begin{proposition}[Calibration exacte de l'échelle]\label{prop-kappa}
Remplacer $\kappa=1$ par $\kappa$ divise chaque $z_g$ par $\sqrt\kappa$ et $\norm{\eh}^2$ par $\kappa$, sans changer
$c_\alpha$. Sur un historique sain, le taux d'alerte est donc une fonction décroissante de $\kappa$, et le plus petit
$\kappa$ dont le taux d'alerte ne dépasse pas la cible $\alpha$ s'obtient exactement par dichotomie.
\end{proposition}
\begin{proof}
Blanchir avec $\kappa\Gamma$ multiplie $W$ par $\kappa^{-1/2}$ ; $\Pi$ et $\Omega$ ne changent pas.
\end{proof}

\subsection{Explication, courbe réparée et erreur de marquage FO}\label{sec-attribution}
\paragraph{Explication.} Après une alerte, les formes types sont retenues une à une : à chaque étape, on ajoute celle
dont le $|z|$ est le plus grand après projection du modèle nul et des formes déjà retenues, tant qu'il dépasse
$c_\alpha$ (poursuite adaptative~\cite{mallat1993}). Les formes retenues sont ensuite réajustées conjointement avec le
modèle nul par moindres carrés généralisés, et chaque amplitude est traduite en unités naturelles
(\cref{tab-templates}).

\paragraph{Courbe réparée.} La courbe Totem réparée retire l'anomalie et garde tout le reste :
\[
q^{\text{rép}}_k=q^{\tot}_k-\frac{1}{T_k}\sum_{g\in\mathcal S}\hat a_g\,g_k .
\]
Les formes contenues dans le sous-espace du modèle nul sont indétectables par construction : un écart de repo lisse
est un désaccord légitime, publié mais jamais signalé comme anomalie de courbe.

\paragraph{Erreur de marquage FO.} Le niveau $\hat\beta_0$ (Totem moins FO, à un an) et son écart-type proviennent du
même ajustement. Le repo FO est signalé comme erreur de marquage possible lorsque $|\hat\beta_0|$ dépasse à la fois
$10$\,pb et $t_{\nu,1-\alpha/2}\,s\,\sigma(\hat\beta_0)$, où :
\begin{itemize}
\item $t_{\nu,1-\alpha/2}$ est un quantile de la loi de Student à $\nu=n^{\opt}-1$ degrés de liberté sur le court
terme : l'écart-type repose sur des dispersions estimées à partir d'une poignée de contributeurs
(\cref{app-student}) ;
\item $s\ge1$ est une échelle calibrée sur l'historique sain comme $\kappa$, pour que les sous-jacents sains déclenchent
au plus une proportion $\alpha$ d'alertes d'erreur de marquage ; elle vaut \vLevelScaleEuStock, \vLevelScaleUsStock,
\vLevelScaleHtb{} et \vLevelScaleIndex{} pour les quatre catégories du banc d'essai.
\end{itemize}
Les algorithmes complets (bootstrap, détection, calibration) sont en \cref{app-algos}.

\subsection{Un exemple à la main}
\begin{example}[Quatre piliers]\label{ex-toy}
Piliers $T=(\num{0.25}\,;\num{0.5}\,;1\,;2)$, bruit indépendant d'écarts-types $(1\,;2\,;4\,;8)$\,pb de portage,
modèle nul réduit au niveau de repo ($X=T$), niveau vrai de $30$\,pb, pic de $12$\,pb en $T_2$ et bruit
$(\num{0.4}\,;\num{-1.1}\,;\num{2.6}\,;\num{-3.0})$\,pb, d'où $\delta=(\num{7.9}\,;\num{25.9}\,;\num{32.6}\,;\num{57.0})$\,pb.
Le niveau ajusté vaut $\hat\beta_0=\vToyBeta$\,pb et $\norm{\eh}^2=\vToyQ$ (le quantile d'ordre \pct{99} de la loi
$\chi^2(3)$ vaut \num{11.3}). Les \vToyM{} formes types donnent :
\begin{center}\input{generated/\lang/tab_toy}\end{center}
Le seuil commun à $\alpha=\pct{2}$ vaut $c_\alpha=\vToyC$ (Šidák : \vToyCsidak). Le pic en $T_2$ l'emporte
($z=\vToyZbest$, amplitude $\vToyAmpBest\pm\vToySeBest$\,pb pour $12$ injectés) ; la marche à partir de $T_2$ obtient
aussi un score élevé, car les deux formes sont corrélées. La poursuite adaptative retient le pic ; ensuite, tous les
$|z|$ restants sont inférieurs à \vToyZafter{} et $\norm{\eh}^2$ tombe à \vToyQafter{} : une seule explication suffit.
\end{example}

\subsection{Retour à l'exemple fil rouge}\label{sec-retour}
\begin{filrouge}
Le \cref{tab-running} donne, pour l'action de l'exemple fil rouge, le diagnostic du détecteur dans les neuf
situations de la \cref{fig-gallery}, avec l'échelle $\kappa=\vRunKappa$ calibrée pour les actions européennes et
$\alpha=\pct{2}$. Le marché sain ne déclenche rien ($\max_g|z_g|=\vRunMaxZ<c_\alpha=\vRunCrit$, et le test global a une
$p$-valeur de \vRunChiP).
\end{filrouge}

\begin{table}[ht]
\centering
\caption{Exemple fil rouge : ce qui est injecté et ce qu'en dit le détecteur. $z$ est la statistique \eqref{eq-z} de
la première forme retenue ; les amplitudes sont traduites en unités naturelles (\cref{tab-templates}).}
\label{tab-running}
\input{generated/\lang/tab_running}
\end{table}

Plusieurs points méritent d'être relevés.
\begin{itemize}
\item \textbf{Les mécanismes sont bien nommés.} Spot, période de dividendes, date de détachement, profil et forward
aberrant sont tous attribués à la bonne forme type et à la bonne maturité. Le décalage du spot est retrouvé à
\vRunSpotEst{} pour \vRunSpotEps{} injectés ; le forward aberrant à \vRunFwdEst\,pb (écart-type \vRunFwdSe) pour
\vRunFwdShift\,pb.
\item \textbf{Les amplitudes sont des estimations.} L'erreur sur 2027 est estimée à \vRunBucketEst{} (écart-type
\vRunBucketSe) pour \vRunBucketDelta{} injectés : le bruit propre du marché sain à neuf mois (l'écart de panel de la
\cref{fig-clean}) se retranche de l'estimation. La part du dividende mal placé est estimée à \vRunTimingEst{} pour
\vRunTimingShare{} injectés, parce que le filtre des cotations aberrantes du consensus a déjà écarté une partie des
contributeurs en désaccord. Le décalage de profil est estimé à \vRunSeasonEst~jours pour \vRunSeasonShift : la
conversion en jours est une linéarisation, juste en signe et en ordre de grandeur.
\item \textbf{Ce qui est lisse n'est pas signalé.} Le sous-panel long biaisé produit une déformation lisse après la
sortie de contributeurs de \vRunPanelTenor{} ; elle est absorbée par l'écart de repo lisse et le détecteur ne signale
rien ($\max_g|z_g|=\vRunPanelMaxZ$). C'est le \cref{thm-nonid} à l'œuvre. De même, le biais brut/net est vu comme une
marche au premier détachement (\vRunGrossEst{} sur le dividende de 2027) plus une pente de repo, comme le prévoit la
\cref{prop-signatures}, (vi).
\item \textbf{Une erreur FO n'est pas une courbe aberrante.} Quand c'est le repo FO qui est faux, aucune anomalie de
courbe n'est signalée, et le test de niveau trouve $\hat\beta_0=\vRunFoLevel$\,pb (écart-type \vRunFoSe), au-delà du
seuil de \vRunFoThr\,pb : erreur de marquage FO possible.
\end{itemize}

\begin{figure}[ht]
\centering\includegraphics[width=\linewidth]{figures/\lang/decomposition.pdf}
\caption{Exemple fil rouge, erreur sur la période 2027 de la \cref{fig-gallery}(b), vue en portage : écart observé,
écart de repo lisse ajusté avec sa bande de $\pm2$ écarts-types, et ajustement lisse plus la marche retenue. Aux
maturités longues, le bruit est si grand que les points s'écartent de l'ajustement sans que cela soit significatif.}
\label{fig-decomposition}
\end{figure}

\begin{retenir}
\begin{itemize}[leftmargin=12pt]
\item Le détecteur juge l'écart de portage avec un bruit calculé à partir des statistiques Totem, retire l'écart de
repo lisse, puis teste une forme type par mécanisme.
\item Chaque test est un filtre adapté exact ; le seuil commun et l'échelle $\kappa$ fixent exactement le taux de
fausses alertes ; la puissance est connue en forme fermée.
\item Chaque alerte vient avec une explication en unités naturelles et une courbe réparée ; le niveau $\hat\beta_0$
sépare les erreurs de marquage FO des anomalies Totem.
\end{itemize}
\end{retenir}
